\documentclass[10pt,twoside]{article}
\usepackage{li_xin_style}
\newcommand{\R}{\mathbb R}
\newcommand{\dd}{\,dx}
\newcommand{\dt}{\,dt}
\newcommand{\Div}{\operatorname{div}}
\newcommand{\supp}{\operatorname{supp}}
\newcommand{\norm}[2]{\lVert#1\rVert_{#2}}
\newcommand{\cH}{\mathcal H}
\renewcommand{\shorttitle}{Research draft: Planar cancellation and energy estimates}
\fancypagestyle{firstpage}{\fancyhf{}
 \fancyhead[L]{\sffamily\fontsize{8}{9.6}\selectfont Research draft}
 \fancyhead[R]{\sffamily\fontsize{8}{9.6}\selectfont\thepage}
 \renewcommand{\headrulewidth}{1pt}}
\title{Planar Cancellation and Energy Estimates\\
for Heat-Conducting Compressible Flow}
\author{}
\date{}
\begin{document}
\maketitle
\begin{abstract}
We study cancellation and energy estimates relevant to the two-dimensional
full compressible Navier--Stokes system with vacuum at infinity.
We prove that an integrable distributional divergence of a square-integrable
vector field on the plane has zero integral, by constructing cutoffs with
scale-invariant gradient energy and controlling their flux tails.
For smooth rapidly decaying solutions of the full system, we derive the
kinetic and internal energy balances, the pressure equation with its thermal
sources, and an estimate for the square of the total specific energy.
These identities exhibit useful cancellations and the terms requiring
further estimates: the pressure-square integral, temperature-weighted viscous
heating, and the velocity-weighted energy flux.
The original local strong-solution class is recorded separately from the
additional hypotheses used to justify these computations.
The estimates do not bound the continuation quantities in that class,
and no small-energy global-existence result is established.
\end{abstract}
\keywords{Compressible Navier--Stokes equations; vacuum; planar cancellation;
heat conduction; energy estimates}

\section{Introduction}
Consider the full compressible Navier--Stokes equations on $\R^2$:
\begin{align}
 \rho_t+\Div(\rho u)&=0, \label{eq:mass}\\
 (\rho u)_t+\Div(\rho u\otimes u)
 -\mu\Delta u-(\mu+\lambda)\nabla\Div u+\nabla P&=0,
 \label{eq:momentum}\\
 c_v\bigl((\rho\theta)_t+\Div(\rho\theta u)\bigr)
 -\kappa\Delta\theta+P\Div u&=Q(u), \label{eq:heat}\\
 P=R\rho\theta,\qquad
 Q(u)&=2\mu|D(u)|^2+\lambda(\Div u)^2. \label{eq:constitutive}
\end{align}
Here $\rho\ge0$, $u\in\R^2$, and $\theta\ge0$ denote density, velocity,
and temperature, respectively, and
\begin{equation}\label{eq:coefficients}
 D(u)=\frac{\nabla u+(\nabla u)^{\mathrm{tr}}}{2},\qquad
 \mu>0,\quad \mu+\lambda\ge0,\quad R,c_v,\kappa>0.
\end{equation}
The convention is $(\nabla u)_{ij}=\partial_j u_i$.
The initial condition is prescribed in conservative variables:
\begin{equation}\label{eq:initial}
 (\rho,\rho u,\rho\theta)(x,0)
   =(\rho_0,\rho_0u_0,\rho_0\theta_0)(x).
\end{equation}
The far-field condition in the local theory of Wang \cite{Wang} is
$(\rho,u,\theta)(x,t)\to(0,0,0)$ as $|x|\to\infty$, $t>0$, with the
qualification \lq\lq in some weak sense\rq\rq.
We retain that qualification; it is not replaced here by pointwise decay
or by an independently chosen normalization of $u$ or $\theta$.

The initial data in \cite[Theorem 1.1]{Wang} satisfy, for $a>1$, $q>2$,
and $\eta_0>0$,
\begin{equation}\label{eq:weight}
 \bar{x}=(e+|x|^2)^{1/2}\log^{1+\eta_0}(e+|x|^2),
\end{equation}
and
\begin{align}
 &0\le\bar{x}^{a}\rho_0\in L^1\cap H^1\cap W^{1,q},
       \qquad \theta_0\ge0, \label{eq:data-density}\\
 &\sqrt{\rho_0}u_0,\ \sqrt{\rho_0}\theta_0\in L^2,\qquad
       \nabla u_0\in H^1,\quad \nabla\theta_0\in L^2,
       \label{eq:data-velocity}\\
 &-\mu\Delta u_0-(\mu+\lambda)\nabla\Div u_0
       +R\nabla(\rho_0\theta_0)=\sqrt{\rho_0}\,g,
       \qquad g\in L^2. \label{eq:compatibility}
\end{align}
All unqualified function spaces below are over $\R^2$.
In particular, \eqref{eq:data-velocity} requires
$\int\rho_0\theta_0^2\dd<\infty$.
It does not separately require $\theta_0\in L^2$ or $u_0\in L^2$.
There is no second thermal compatibility condition in these printed
hypotheses.

Put
\begin{equation}\label{eq:exponents}
 \alpha=\min\{a/2,1\},\qquad p_u=\frac4\alpha,\qquad
 p_\theta=\frac6{2\alpha-1},\qquad q_1\in(2,\infty).
\end{equation}
Thus $1/2<\alpha\le1$ and both $p_u,p_\theta$ are finite.
On a local existence interval $[0,T_0]$, the class stated in
\cite[(1.8)--(1.9)]{Wang} is
\begin{align}
 &\rho\in C([0,T_0];L^1\cap H^1\cap W^{1,q}),\nonumber\\
 &\bar{x}^{a}\rho\in L^\infty(0,T_0;L^1\cap H^1\cap W^{1,q}),
       \label{eq:class-density}\\
 &\nabla u\in L^\infty(0,T_0;H^1)
                  \cap L^2(0,T_0;W^{1,q}),\qquad
 u\in L^\infty(0,T_0;L^{p_u}),\label{eq:class-velocity}\\
 &\sqrt\rho u,\ \sqrt\rho\theta,\ \sqrt\rho u_t,\ \nabla\theta,\
   \sqrt t\sqrt\rho\theta_t,\ \sqrt t\nabla^2\theta,\
   \bar{x}^{\alpha/2}\nabla u
      \in L^\infty(0,T_0;L^2), \label{eq:class-bounded}\\
 &\nabla u_t,\ \sqrt\rho\theta_t,\ \nabla^2\theta,\
       \sqrt t\nabla\theta_t
      \in L^2(\R^2\times(0,T_0)), \label{eq:class-integrated}\\
 &\theta\in L^2(0,T_0;L^{p_\theta}),\qquad
   \sqrt t\nabla^2\theta\in L^2(0,T_0;L^{q_1}).
       \label{eq:class-temperature}
\end{align}
The strong-solution definition in that source requires the derivatives in
the equations to be regular distributions and the equations to hold almost
everywhere. The condition in \eqref{eq:class-velocity} concerns $u$
itself, not a class of velocities modulo constants.

Let $T_{\max}$ denote the maximal existence time in this class.
The continuation criterion stated in \cite[Theorem 1.2]{Wang} is
\begin{equation}\label{eq:criterion}
 T_{\max}<\infty\quad\Longrightarrow\quad
 \lim_{T\uparrow T_{\max}}\int_0^T
       \bigl(\norm{\Div u}{\infty}+\norm{u}{p_u}\bigr)\dt=\infty.
\end{equation}
Both time exponents in this criterion are one.
A bound on compact subintervals whose constants may diverge as
$T\uparrow T_{\max}$ does not establish the required continuation bound.

The small quantity motivating this investigation is the physical
initial energy
\begin{equation}\label{eq:E0}
 E_0=\int_{\R^2}
       \left(\frac12\rho_0|u_0|^2+c_v\rho_0\theta_0\right)\dd.
\end{equation}
It is finite under \eqref{eq:data-density}--\eqref{eq:data-velocity}:
with $m_0=\int\rho_0\dd$, Cauchy--Schwarz gives
$\int\rho_0\theta_0\dd
 \le m_0^{1/2}\norm{\sqrt{\rho_0}\theta_0}{2}$.
The question is whether sufficiently small $E_0$, for prescribed finite
bounds on the other initial quantities, ensures global continuation in
\eqref{eq:class-density}--\eqref{eq:class-temperature}.
No small higher norm, small mass, zero total momentum, second moment,
positive density lower bound, or support restriction is appended to that
question. No mass normalization is used below.
The printed initial hypotheses in \cite{Wang} do not separately state
$m_0>0$; an argument using a positive local mass must specify that premise
and the dependence of its constants. Our cancellation theorem uses no
mass condition.

The planar cutoff argument behind the cancellation methods of
Fan--L\"u--Wang \cite{FLW} concerns integrable divergences of
square-integrable fields. Section~\ref{sec:cancellation} proves that
general result without a total-momentum assumption.
The energy approach, for which Li--Xin \cite{LiXin} is a useful reference
in the compressible setting, must retain pressure work, viscous heating,
and thermal conduction in \eqref{eq:heat}.
Sections~\ref{sec:energy}--\ref{sec:specific} derive corresponding
estimates with their hypotheses displayed. They do not establish the
missing uniform bounds in \eqref{eq:criterion}.

Throughout, $\norm{\cdot}{p}$ denotes the $L^p(\R^2)$ norm, and
$\dot f=f_t+u\cdot\nabla f$ denotes the material derivative.
Unless otherwise specified, constants $C$ depend only on the physical
coefficients in \eqref{eq:coefficients}; $C_\varepsilon$ may additionally
depend on the positive parameter $\varepsilon$.

\section{A planar cancellation theorem}\label{sec:cancellation}
We first isolate the limiting argument. Its hypotheses concern actual
Lebesgue integrals and do not include a differential equation.

\begin{lemma}\label{lem:tail}
Let $f,h\in L^1(\R^2)$ be real-valued, with $h\ge0$ almost everywhere.
Let $\chi_n$ be measurable real functions such that
$|\chi_n(x)|\le B$ for every $n\in\mathbb N$ and $x\in\R^2$,
for a fixed $B\ge0$, and $\chi_n\to1$ almost everywhere.
Suppose that, for a fixed $C\ge0$,
\begin{equation}\label{eq:tail-hypothesis}
 \left|\int f\chi_n\dd\right|
   \le C\left(\int_{\{|x|\ge n\}}h\dd\right)^{1/2}
       \qquad(n\in\mathbb N).
\end{equation}
Then $\int f\dd=0$.
\end{lemma}
\begin{proof}
The products $f\chi_n$ are integrable because $|f\chi_n|\le B|f|$.
For each fixed $x$, the indicator of $\{|x|\ge n\}$ tends to zero.
The functions $h\mathbf1_{\{|x|\ge n\}}$ are dominated by $h$,
so dominated convergence gives
\[
 \lim_{n\to\infty}\int_{\{|x|\ge n\}}h\dd=0.
\]
These integrals are nonnegative and their square roots also tend to zero.
On the other hand, dominated convergence with dominating function
$B|f|$ yields
\[
 \lim_{n\to\infty}\int f\chi_n\dd=\int f\dd.
\]
Taking limits in \eqref{eq:tail-hypothesis} proves the conclusion.
\end{proof}

\begin{theorem}\label{thm:cancellation}
Let $f\in L^1(\R^2;\R)$ and $V\in L^2(\R^2;\R^2)$.
Suppose that $f=\Div V$ in distributions, that is,
\begin{equation}\label{eq:distributional}
 \int f\varphi\dd=-\int V\cdot\nabla\varphi\dd
       \qquad\text{for every }\varphi\in C_c^\infty(\R^2).
\end{equation}
Then $\int_{\R^2}f\dd=0$.
\end{theorem}
\begin{proof}
Construct a cutoff by setting
\[
 b(t)=
 \begin{cases}\exp(-1/t),&t>0,\\0,&t\le0,\end{cases}
 \qquad
 \eta(x)=\frac{b(4-|x|^2)}
                 {b(4-|x|^2)+b(|x|^2-1)}.
\]
Every derivative of $b$ from the right vanishes at zero, so $b$ is smooth.
The two arguments in the denominator cannot both be nonpositive.
Thus $\eta$ is smooth, $0\le\eta\le1$, $\eta=1$ on $|x|\le1$,
and $\eta=0$ on $|x|\ge2$. In particular,
$K_\eta=\int|\nabla\eta|^2\dd<\infty$.
For $L>0$, define $\chi_L(x)=\eta(x/L)$.
The chain rule gives
\begin{equation}\label{eq:cutoff-gradient}
 \nabla\chi_L(x)=L^{-1}\nabla\eta(x/L),\qquad
 \supp\nabla\chi_L\subset\{L\le|x|\le2L\}.
\end{equation}
The planar change of variables $x=Ly$ yields
\begin{equation}\label{eq:cutoff-scaling}
 \int|\nabla\chi_L(x)|^2\dd
 =L^{-2}L^2\int|\nabla\eta(y)|^2\,dy=K_\eta.
\end{equation}
All integrals here are finite by smoothness and compact support.

Take $L=n+1$. The actual test function $\chi_{n+1}$ is admissible in
\eqref{eq:distributional}, so
\begin{align}
 \left|\int f\chi_{n+1}\dd\right|
 &\le\int_{\{|x|\ge n\}}|V|\,|\nabla\chi_{n+1}|\dd\nonumber\\
 &\le K_\eta^{1/2}
       \left(\int_{\{|x|\ge n\}}|V|^2\dd\right)^{1/2}.
       \label{eq:derived-tail}
\end{align}
The tested forcing is integrable since $|\chi_{n+1}|\le1$.
The flux pairing is integrable by Cauchy--Schwarz, using
$V\in L^2$ and \eqref{eq:cutoff-scaling}.
Moreover, $\chi_{n+1}(x)=1$ for all sufficiently large $n$ at every
fixed $x$. Lemma~\ref{lem:tail}, with
$h=|V|^2$, $B=1$, and $C=K_\eta^{1/2}$, proves the conclusion.
\end{proof}

The scale-independent bound \eqref{eq:cutoff-scaling} is the planar
feature of this argument. In dimension $d$, the same dilation gives
$L^{d-2}K_\eta$. No $L^1$ assumption on $V$ was used in the plane,
and the constant function one was never used as a distributional test.

Theorem~\ref{thm:cancellation} is a statement about $f$ and $V$ with
the displayed function-space hypotheses. It is not an assertion that
a particular forcing or stress of \eqref{eq:mass}--\eqref{eq:heat}
satisfies them. Such an application requires its own verification of
the distributional identity, $L^1$ integrability of the forcing, and
$L^2$ integrability of the flux. The theorem alone supplies neither
weighted elliptic bounds nor the norms in \eqref{eq:criterion}.

\section{Physical energy and pressure work}\label{sec:energy}
For Sections~\ref{sec:energy}--\ref{sec:specific}, impose the following
sufficient regularity condition explicitly: $(\rho,u,\theta)$ is a smooth
solution with $\rho,\theta\ge0$, and these functions and all spatial
and time derivatives used below decay faster than every inverse power
of $|x|$, uniformly on compact time intervals.
This condition makes every displayed integral finite and permits
differentiation under the integral. Its validity for the original class
\eqref{eq:class-density}--\eqref{eq:class-temperature} is not asserted.
The conditional identities below assert no existence of positive-energy
solutions in this auxiliary decay class. They provide neither a local
construction nor a non-vacuity assertion for that class.

Integrations on $\R^2$ can then be justified with $\chi_L$ from
\eqref{eq:cutoff-gradient}. Every smooth rapidly decaying vector field
$F$ used below satisfies
\begin{equation}\label{eq:vanishing-flux}
 \begin{split}
 \int\chi_L\Div F\dd&=-\int F\cdot\nabla\chi_L\dd,\\
 \left|\int F\cdot\nabla\chi_L\dd\right|
 &\le\frac{\norm{\nabla\eta}{\infty}}{L}
          \int_{\{|x|\ge L\}}|F|\dd\longrightarrow0.
 \end{split}
\end{equation}
Dominated convergence treats the volume integrals.
For integration by parts of two factors, apply this estimate to their
product flux. Uniform decay supplies time-local dominating functions.
These observations justify the boundary and time operations used below.

Write
\begin{equation}\label{eq:notation-stress}
 d=\Div u,\quad k=\tfrac12|u|^2,\quad
 S=2\mu D(u)+\lambda dI,\quad \nu=2\mu+\lambda>0.
\end{equation}
Then $\Div S=\mu\Delta u+(\mu+\lambda)\nabla d$ and $Q=S:\nabla u$.
In two dimensions,
\begin{equation}\label{eq:Q-positive}
 Q=2\mu\left|D(u)-\tfrac12dI\right|^2+(\mu+\lambda)d^2\ge0.
\end{equation}
Define
\begin{align}
 K(t)&=\int\rho k\dd,\qquad H(t)=c_v\int\rho\theta\dd,\qquad
 E(t)=K(t)+H(t),\label{eq:energy-quantities}\\
 \mathcal D(t)&=\int\{\mu|\nabla u|^2+(\mu+\lambda)d^2\}\dd.
 \label{eq:dissipation}
\end{align}

\begin{proposition}\label{prop:physical-energy}
Under the regularity condition above,
\begin{equation}\label{eq:energy-balances}
 K'+\mathcal D=\int Pd\dd,\qquad
 H'=\mathcal D-\int Pd\dd,\qquad E'=0.
\end{equation}
If the interval includes zero, with the stated initial traces, then
\begin{equation}\label{eq:energy-small}
 K(t)+H(t)=E_0,\qquad
 \norm{\sqrt\rho u}{2}^2\le2E_0,\qquad
 \norm{P}{1}\le\frac{R}{c_v}E_0.
\end{equation}
For any two times $s<t$ in the interval,
\begin{equation}\label{eq:kinetic-pressure}
 K(t)+\frac{\mu}{2}\int_s^t\norm{\nabla u}{2}^2\dt
 \le K(s)+\frac1\mu\int_s^t\norm{P}{2}^2\dt.
\end{equation}
\end{proposition}
\begin{proof}
Continuity gives the transport identity
\begin{equation}\label{eq:transport}
 \partial_t(\rho f)+\Div(\rho uf)=\rho\dot f.
\end{equation}
The momentum equation is $\rho\dot u=\Div S-\nabla P$.
Dotting it with $u$, and using the symmetry of $S$, gives
\begin{equation}\label{eq:kinetic-local}
 \rho\dot k=\Div(Su-Pu)-Q+Pd.
\end{equation}
Integrate \eqref{eq:transport} with $f=k$ and apply
\eqref{eq:vanishing-flux} to obtain $K'+\int Q\dd=\int Pd\dd$.
Integration by parts twice gives
\[
 \int\partial_i u_j\,\partial_j u_i\dd
 =-\int u_j\,\partial_j(\Div u)\dd=\int d^2\dd,
\]
where repeated indices are summed over $1,2$. Since
$2|D(u)|^2=|\nabla u|^2+\sum_{i,j}\partial_i u_j\,\partial_j u_i$,
\begin{equation}\label{eq:integrated-Q}
 \int Q\dd=\mu\norm{\nabla u}{2}^2
                 +(\mu+\lambda)\norm{d}{2}^2=\mathcal D.
\end{equation}
The heat equation, \eqref{eq:transport} with $f=\theta$, and
\eqref{eq:vanishing-flux} with $F=\nabla\theta$ give
$H'=\int Q\dd-\int Pd\dd$. This proves \eqref{eq:energy-balances}.
Add the first two identities, integrate in time, and use $\rho,\theta\ge0$
to obtain \eqref{eq:energy-small}.

Finally, $|d|\le\sqrt2|\nabla u|$, and Cauchy--Schwarz followed by
Young's inequality gives
\[
 \left|\int Pd\dd\right|
 \le\sqrt2\norm{P}{2}\norm{\nabla u}{2}
 \le\frac{\mu}{2}\norm{\nabla u}{2}^2+\frac1\mu\norm{P}{2}^2.
\]
Since $\mathcal D\ge\mu\norm{\nabla u}{2}^2$, integration of the
kinetic identity proves \eqref{eq:kinetic-pressure}.
\end{proof}

The physical energy conserves the sum of kinetic and internal energy.
The term $\mathcal D$ transfers energy between them; it does not
appear with a positive sign in the equation for $E$.
Thus \eqref{eq:energy-small} supplies no bound on
$\int_s^t\mathcal D\dt$ by itself. If $M(t)=\norm{\rho(t)}{\infty}$,
\begin{equation}\label{eq:pressure-thermal}
 \norm{P}{2}^2\le R^2M(t)\int\rho\theta^2\dd,
\end{equation}
but this introduces a higher thermal quantity and a density coefficient,
neither bounded uniformly up to $T_{\max}$ by \eqref{eq:energy-small}.

The effective viscous flux $G=\nu d-P$ describes the same coupling.

\begin{proposition}\label{prop:pressure-feedback}
Under the same regularity condition,
\begin{align}
 H'+\frac1\nu\norm{P}{2}^2
     &=\mathcal D-\frac1\nu\int PG\dd,\label{eq:internal-flux}\\
 K'+\frac12\mathcal D
     &\le\frac1{2\nu}\norm{P}{2}^2.\label{eq:kinetic-flux}
\end{align}
Consequently,
\begin{equation}\label{eq:pressure-spacetime}
 \frac1{2\nu}\int_s^t\norm{P}{2}^2\dt
 \le H(s)+\int_s^t\mathcal D\dt
                 +\frac1{2\nu}\int_s^t\norm{G}{2}^2\dt.
\end{equation}
\end{proposition}
\begin{proof}
Substitute $d=(G+P)/\nu$ in the internal-energy identity to obtain
\eqref{eq:internal-flux}. The identity
\[
 \norm{\nabla u}{2}^2=\norm{d}{2}^2+\norm{\omega}{2}^2,\qquad
 \omega=\partial_1u_2-\partial_2u_1,
\]
follows by expanding the squares and using
$\int\partial_1u_1\,\partial_2u_2\dd
 =\int\partial_2u_1\,\partial_1u_2\dd$.
Hence $\mathcal D=\nu\norm{d}{2}^2+\mu\norm{\omega}{2}^2
\ge\nu\norm{d}{2}^2$, and
\[
 \int Pd\dd\le\frac1{2\nu}\norm{P}{2}^2+\frac{\nu}{2}\norm{d}{2}^2
 \le\frac1{2\nu}\norm{P}{2}^2+\frac12\mathcal D.
\]
This proves \eqref{eq:kinetic-flux}. Finally,
$-\int PG\dd\le\frac12\norm{P}{2}^2+\frac12\norm{G}{2}^2$.
Insert this in \eqref{eq:internal-flux}, integrate, and discard the
nonnegative $H(t)$ to obtain \eqref{eq:pressure-spacetime}.
\end{proof}

Integration of \eqref{eq:kinetic-flux} yields
\[
 \int_s^t\mathcal D\dt
 \le2K(s)+\frac1\nu\int_s^t\norm{P}{2}^2\dt.
\]
Inserting this in \eqref{eq:pressure-spacetime} leaves the coefficient
$1/\nu$ of the pressure integral on the right, exceeding its coefficient
$1/(2\nu)$ on the left. The two estimates therefore give no absorption.
At the level of exact identities their sum is $E'=0$.
A further estimate independent of this exchange is needed.

\section{Thermal sources and the pressure equation}\label{sec:thermal}
Set $\beta=R/c_v$. Since $\dot\rho=-\rho d$ and
$c_v\rho\dot\theta=\kappa\Delta\theta+Q-Pd$, the pressure equation is
\begin{equation}\label{eq:pressure-evolution}
 \dot P+(1+\beta)Pd=\beta(Q+\kappa\Delta\theta).
\end{equation}
No division by density is involved. The two thermal sources cannot be
discarded in favor of a source-free barotropic pressure equation.

\begin{proposition}\label{prop:pressure-square}
Under the regularity condition of Section~\ref{sec:energy},
\begin{align}
 \frac d{dt}\norm{P}{2}^2+\frac{1+2\beta}{\nu}\int P^3\dd
 ={}&-\frac{1+2\beta}{\nu}\int P^2G\dd+2\beta\int PQ\dd\nonumber\\
 &+2\beta\kappa\int P\Delta\theta\dd,\label{eq:pressure-square}\\
 \int P\Delta\theta\dd
 ={}&-R\int\rho|\nabla\theta|^2\dd
       -R\int\theta\nabla\rho\cdot\nabla\theta\dd.
       \label{eq:conduction-pressure}
\end{align}
\end{proposition}
\begin{proof}
Multiply \eqref{eq:pressure-evolution} by $2P$ and use
$\int u\cdot\nabla(P^2)\dd=-\int dP^2\dd$. Then
\[
 \frac d{dt}\norm{P}{2}^2+(1+2\beta)\int P^2d\dd
       =2\beta\int PQ\dd+2\beta\kappa\int P\Delta\theta\dd.
\]
Substitute $d=(G+P)/\nu$ to obtain \eqref{eq:pressure-square}.
Integration by parts with
$\nabla P=R\rho\nabla\theta+R\theta\nabla\rho$ proves
\eqref{eq:conduction-pressure}. All product fluxes satisfy
\eqref{eq:vanishing-flux}.
\end{proof}

The favorable cubic term in \eqref{eq:pressure-square} is accompanied
by $\int PQ\dd\ge0$ on the right. The last integral in
\eqref{eq:conduction-pressure} has no fixed sign.
Neither term is controlled by \eqref{eq:energy-small}.

For the thermal estimate, write
\begin{equation}\label{eq:thermal-notation}
 \begin{gathered}
 X=\int\rho\theta^2\dd,\quad Z=\norm{\nabla\theta}{2}^2,\quad
 J=\int\theta Q\dd,\\
 V=\int\rho|\dot u|^2\dd,\qquad N=\norm{\nabla u}{2}^2.
 \end{gathered}
\end{equation}
\begin{proposition}\label{prop:thermal}
Under the regularity condition of Section~\ref{sec:energy},
\begin{equation}\label{eq:thermal-identity}
 \frac{c_v}{2}X'+\kappa Z=J-R\int\rho\theta^2d\dd.
\end{equation}
Writing $U=\norm{u}{\infty}$ and $M=\norm{\rho}{\infty}$, for every
$\varepsilon>0$ one has
\begin{equation}\label{eq:heating-bound}
 J\le\varepsilon(V+Z)+C_\varepsilon U^2\{N+(1+M)X\}
                  +C\norm{d}{\infty} X.
\end{equation}
\end{proposition}
\begin{proof}
Multiply the heat equation by $\theta$ and apply
\eqref{eq:transport} to $\theta^2/2$. Integration by parts of conduction
gives \eqref{eq:thermal-identity}.
Testing the momentum equation by $u\theta$ gives
\begin{align}
 \int\theta\{\mu|\nabla u|^2+(\mu+\lambda)d^2\}\dd
 ={}&-\int\rho\theta u\cdot\dot u\dd+\int P\theta d\dd
                 +\int Pu\cdot\nabla\theta\dd\nonumber\\
 &-\mu\int\partial_i u_j\,u_j\partial_i\theta\dd
       -(\mu+\lambda)\int d\,u\cdot\nabla\theta\dd.
       \label{eq:momentum-thermal}
\end{align}
Because $\theta\ge0$ and $\mu+\lambda\ge0$, the left side is at least
$\mu\int\theta|\nabla u|^2\dd$. Also,
$0\le Q\le(2\mu+2|\lambda|)|\nabla u|^2$.
Thus $J$ is bounded by a constant times the sum of the absolute values
on the right of \eqref{eq:momentum-thermal}. Cauchy--Schwarz gives
\begin{align*}
 \left|\int\rho\theta u\cdot\dot u\dd\right|
   &\le UX^{1/2}V^{1/2},\\
 \left|\int P\theta d\dd\right|&\le R\norm{d}{\infty} X,\\
 \left|\int Pu\cdot\nabla\theta\dd\right|
   &\le RU(MX)^{1/2}Z^{1/2}.
\end{align*}
The two viscous cross terms are bounded by $CUN^{1/2}Z^{1/2}$,
using $|d|\le\sqrt2|\nabla u|$.
Young's inequality gives
\[
 CUX^{1/2}V^{1/2}\le\varepsilon V+C_\varepsilon U^2X
\]
and, choosing the two parameters with sum $\varepsilon$,
\[
 CU(MX)^{1/2}Z^{1/2}+CUN^{1/2}Z^{1/2}
 \le\varepsilon Z+C_\varepsilon U^2(MX+N).
\]
Together these inequalities prove \eqref{eq:heating-bound}.
\end{proof}

With $d_-=\max\{-d,0\}$,
\begin{equation}\label{eq:compression}
 -R\int\rho\theta^2d\dd\le R\norm{d_-}{\infty} X.
\end{equation}
Using this coefficient in Gronwall's inequality requires its time
integrability, already part of \eqref{eq:criterion}.
The quantities $V$, $U$, and $M$ in \eqref{eq:heating-bound} also remain
to be controlled uniformly up to a finite $T_{\max}$.

For a time-independent $h\in C_c^\infty(\R^2)$ the localized
temperature test is
\begin{align}
 \frac{c_v}{2}\frac d{dt}\int h\rho\theta^2\dd
       +\kappa\int h|\nabla\theta|^2\dd
 ={}&\int h\theta Q\dd-\int hP\theta d\dd\nonumber\\
 &+\frac{c_v}{2}\int\rho\theta^2u\cdot\nabla h\dd
       +\frac{\kappa}{2}\int\theta^2\Delta h\dd.
       \label{eq:localized-temperature}
\end{align}
Indeed, apply \eqref{eq:transport} to $\theta^2/2$ and multiply by $h$;
for conduction use
\[
 \kappa\int h\theta\Delta\theta\dd
 =-\kappa\int h|\nabla\theta|^2\dd
       +\frac{\kappa}{2}\int\theta^2\Delta h\dd.
\]
An extension to a growing weight requires control of both additional
terms in \eqref{eq:localized-temperature}.
The density-weighted physical energy gives no unweighted $\theta^2$
control in a vacuum region.

\section{A total specific-energy multiplier}\label{sec:specific}
Pressure work and viscous heating can be cancelled before estimating
them. Set
\begin{equation}\label{eq:specific-definitions}
 w=c_v\theta+k,\quad F=Su-Pu,\quad a_\kappa=\frac{\kappa}{c_v},
 \quad Y=\int\rho w^2\dd.
\end{equation}
\begin{proposition}\label{prop:specific}
Under the regularity condition of Section~\ref{sec:energy},
\begin{equation}\label{eq:specific-equation}
 \rho\dot w=\kappa\Delta\theta+\Div F.
\end{equation}
The functional $Y$ satisfies
\begin{align}
 \frac12Y'+\kappa c_v\norm{\nabla\theta}{2}^2
 &=-\kappa\int\nabla k\cdot\nabla\theta\dd
       -\int F\cdot(c_v\nabla\theta+\nabla k)\dd,
       \label{eq:specific-theta}\\
 \frac12Y'+a_\kappa\norm{\nabla w}{2}^2
 &=\int(a_\kappa\nabla k-F)\cdot\nabla w\dd.
       \label{eq:specific-w}
\end{align}
In particular,
\begin{equation}\label{eq:specific-bound}
 \begin{split}
 Y'+a_\kappa\norm{\nabla w}{2}^2
 &\le a_\kappa^{-1}\norm{a_\kappa\nabla k-F}{2}^2\\
 &\le C\int|u|^2(|\nabla u|^2+P^2)\dd .
 \end{split}
\end{equation}
With $N$ as in \eqref{eq:thermal-notation},
$U=\norm{u}{\infty}$, and $M=\norm{\rho}{\infty}$,
\begin{equation}\label{eq:specific-gronwall}
 \begin{split}
 Y(t)+a_\kappa\int_s^t\norm{\nabla w}{2}^2\dt
 \le{}&\left(Y(s)+C\int_s^tU^2N\dt\right)\\
 &{}\times\exp\left(C\int_s^tMU^2\dt\right).
 \end{split}
\end{equation}
\end{proposition}
\begin{proof}
Add \eqref{eq:kinetic-local} to the material form of the heat equation.
The terms $Q-Pd$ cancel and give \eqref{eq:specific-equation}
without division by density. Since
$Su=\mu\nabla k+\mu(u\cdot\nabla)u+\lambda ud$, equivalently
\[
 \rho\dot w=\kappa\Delta\theta+\mu\Delta k
       +\mu\Div((u\cdot\nabla)u)+\lambda\Div(ud)-\Div(Pu).
\]
Apply \eqref{eq:transport} to $w^2/2$ and test
\eqref{eq:specific-equation} by $w$. Then
\begin{equation}\label{eq:specific-test}
 \frac12Y'=-\kappa\int\nabla w\cdot\nabla\theta\dd
                     -\int F\cdot\nabla w\dd.
\end{equation}
The substitutions $\nabla w=c_v\nabla\theta+\nabla k$ and
$\nabla\theta=c_v^{-1}(\nabla w-\nabla k)$ respectively prove
\eqref{eq:specific-theta} and \eqref{eq:specific-w}.

Cauchy--Schwarz and Young's inequality yield
\[
 \left|\int(a_\kappa\nabla k-F)\cdot\nabla w\dd\right|
 \le\frac{a_\kappa}{2}\norm{\nabla w}{2}^2
       +\frac1{2a_\kappa}\norm{a_\kappa\nabla k-F}{2}^2.
\]
Insert this in \eqref{eq:specific-w} and multiply by two.
The pointwise bounds $|\nabla k|\le|u||\nabla u|$ and
$|S|\le C|\nabla u|$ imply
\[
 |a_\kappa\nabla k-F|^2\le C|u|^2(|\nabla u|^2+P^2),
\]
which proves \eqref{eq:specific-bound}. Positivity also gives
\begin{equation}\label{eq:specific-positive}
 Y\ge c_v^2X+\frac14\int\rho|u|^4\dd,\qquad
 \norm{P}{2}^2\le\frac{R^2}{c_v^2}MY.
\end{equation}
Consequently,
\begin{equation}\label{eq:specific-coefficients}
 Y'+a_\kappa\norm{\nabla w}{2}^2\le CU^2N+CMU^2Y.
\end{equation}
Multiply by $\exp(-C\int_s^tMU^2\,d\tau)$ and integrate.
After multiplying back, the dissipation integral has the weight
$\exp(C\int_\tau^tMU^2\,d\sigma)\ge1$.
The forcing term is at most
$\exp(C\int_s^tMU^2\,d\tau)\int_s^tCU^2N\,d\tau$.
Dropping the excess dissipation weight proves
\eqref{eq:specific-gronwall}.
\end{proof}

The positive functional $Y$ contains both the density-weighted
temperature square and velocity fourth power.
However, $\nabla w=c_v\nabla\theta+\nabla k$ controls only their sum.
The identity \eqref{eq:specific-theta} retains the mixed conduction
term, whereas \eqref{eq:specific-w} places the other terms in a flux.
No small factor depending on $E_0$ has been established on the right
of \eqref{eq:specific-bound}. The coefficients in
\eqref{eq:specific-gronwall} have not been bounded uniformly at a
finite maximal time under the original data assumptions.

Small $\int\rho_0w_0\dd=E_0$ does not imply small
$\int\rho_0w_0^2\dd$. To see the distinction, choose a smooth
nonnegative compactly supported density equal to one near the origin
and a nonnegative nonzero $\psi\in C_c^\infty(\R^2)$.
For $r>0$ sufficiently small, set
\[
 u_0=0,\qquad \theta_0(x)=b\psi(x/r),\qquad b>0,
\]
with the thermal support inside the region where $\rho_0=1$.
The data satisfy \eqref{eq:data-density}--\eqref{eq:data-velocity}.
Condition \eqref{eq:compatibility} holds with $g=R\nabla\theta_0$,
extended by zero away from its support.
A change of variables gives
\begin{equation}\label{eq:concentration}
 \begin{gathered}
 E_0=c_vbr^2\int\psi\dd,\qquad
 \norm{P_0}{2}^2=R^2b^2r^2\int\psi^2\dd,\\
 Y(0)=c_v^2b^2r^2\int\psi^2\dd.
 \end{gathered}
\end{equation}
For fixed positive small $E_0$, take
$r^2=E_0/(c_vb\int\psi\dd)$ and let $b$ increase.
Both $\norm{P_0}{2}^2$ and $Y(0)$ then grow linearly in $b$.
The higher initial norms are finite for each datum but not uniform in
this family. This rules out an instantaneous bound using $E_0$ alone;
it does not rule out a spacetime estimate using smoothing and prescribed
higher initial bounds, and gives no conclusion against global existence.

\section{Remaining estimates and the continuation question}
Theorem~\ref{thm:cancellation} is complete under its $L^1$ forcing and
$L^2$ flux assumptions. The energy and thermal estimates are proved under the explicit
smoothness and decay condition of Section~\ref{sec:energy}. Their use for every solution in
\eqref{eq:class-density}--\eqref{eq:class-temperature} requires the
corresponding cutoff limits, nonlinear tests, and initial traces.
For example, the $w^2$ test requires integrability of $\rho|u|^4$
and its flux terms. A formal multiplier or the weak far-field phrase
does not prove these facts.

Even after those justifications, a chain of unresolved estimates remains.
The kinetic inequality \eqref{eq:kinetic-pressure} requires the
pressure-square time integral. The effective-flux identities
\eqref{eq:internal-flux}--\eqref{eq:pressure-spacetime} do not provide
an absorbable bound for it. The pressure equation introduces
$\int PQ\dd$ and the mixed density-gradient term in
\eqref{eq:conduction-pressure}. The thermal test requires
temperature-weighted heating and a compression coefficient.
The total specific-energy functional replaces the explicit heating
and pressure work by $\int|u|^2(|\nabla u|^2+P^2)\dd$,
which remains uncontrolled by the physical energy estimate.

Continuation requires constants remaining finite as
$T\uparrow T_{\max}<\infty$, sufficient to prove
\[
 \sup_{T<T_{\max}}\int_0^T
       \bigl(\norm{\Div u}{\infty}+\norm{u}{4/\alpha}\bigr)\dt<\infty.
\]
Assuming this integral finite to estimate the preceding coefficients
cannot establish its finiteness from small initial energy.
A bootstrap density bound or velocity supremum bound must likewise
be recovered from the equations with uniform constants.

The results therefore identify positive cancellations and exact
requirements for a future closure. They establish neither a
small-energy global-existence theorem in the original local class
nor a contradiction to such a theorem.
An independent estimate connecting the pressure and thermal flux terms
to the continuation quantities is still needed, without adding
assumptions to the target data.

\begin{thebibliography}{9}
\bibitem{Wang}
X. Wang,
\emph{On Local Existence and Blowup Criterion of Strong Solutions to
Cauthy Problem of 2D Full Compressible Navier-Stokes System with Vacuum},
arXiv:2212.13343v1, 27 December 2022.
\bibitem{LiXin}
J. Li and Z. Xin,
\emph{Global Well-Posedness and Large Time Asymptotic Behavior of Classical
Solutions to the Compressible Navier--Stokes Equations with Vacuum},
Annals of PDE \textbf{5} (2019), Article 7.
\href{https://doi.org/10.1007/s40818-019-0064-5}
{doi:10.1007/s40818-019-0064-5}.
\bibitem{FLW}
X. Fan, B. L\"u, and X. Wang,
\emph{Hardy Space and Exponential Decay of Inhomogeneous Navier--Stokes
System in $\R^2$},
arXiv:2610.04884v1, 4 October 2026.
\end{thebibliography}
\end{document}
